\subsubsection{Plot3D}
\label{SEC:Files:Grid_Input:Plot3d}

%%%%%%%%%%%%%%%%%%%%%%%%%%%%%%%%%%%%%%%%%%%%%%%%%%%%%%%%%%%%%%%%%%%%%%%%%%%%%%%%

When using Plot3D grid files, additional controls are available to
individually strip points in each of the computational directions.
%
The technique is often referred to as adjusting the grid stride or
grid sequencing.
%
The resulting coarser mesh appears as simple hexahedral elements
formed from the remaining grid nodes and the omitted nodes
provide no contribution.
%
For a higher-order method like flux reconstruction, this can have
an adverse effect along boundaries, where the discrete representation
of the surface curvature becomes reduced.
%
The coarser faceting of the geometry surface can introduce
spurious waves into the solution.

After completing the optional coarsening of the input grid, the
flow solver can be instructed to combine a certain multiple of the
remaining adjacent cells in all directions to form agglomerated cells.
%
These agglomerated cells are not simple hexes with planar faces.
%
Rather the faces are formed by polynomials defined by the nodes of the
intermediate cells that were combined together.
%
This is a somewhat crude way of constructing higher-order cells that
preserve surface curvature.
%
In order for agglomeration to work, the number of nodes
in each direction of the grid must be an appropriate multiple of
the desired agglomeration order.
%
For example, a third-order agglomeration requires merging three
cells in each direction.

\NOTE
The agglomeration order for the geometry is independent of the
solution order.
%
The volume solution points and face flux points will be interpolated
from the curved geometry information.

%%%%%%%%%%%%%%%%%%%%%%%%%%%%%%%%%%%%%%%%%%%%%%%%%%%%%%%%%%%%%%%%%%%%%%%%%%%%%%%%

\SeeAlso{
\inp[nmlGrid::Plot3dAgglomerateOrder]{Plot3d\_Agglomerate\_Order},
\inp[nmlGrid::Plot3dStrideAll]{Plot3d\_Stride\_All},
\inp[nmlGrid::Plot3dStrideI]{Plot3d\_Stride\_I},
\inp[nmlGrid::Plot3dStrideJ]{Plot3d\_Stride\_J},
\inp[nmlGrid::Plot3dStrideK]{Plot3d\_Stride\_K},
\inp[nmlNumerics::SolutionOrder]{Solution\_Order}
}

%%%%%%%%%%%%%%%%%%%%%%%%%%%%%%%%%%%%%%%%%%%%%%%%%%%%%%%%%%%%%%%%%%%%%%%%%%%%%%%%


\paragraph{Plot3D Interface Cut File}

%%%%%%%%%%%%%%%%%%%%%%%%%%%%%%%%%%%%%%%%%%%%%%%%%%%%%%%%%%%%%%%%%%%%%%%%%%%%%%%%

For multiblock Plot3D grids, block connectivity is specified via a
text-based interface cut file.
%
An example of this cut file is shown in \lref{LST:Plot3d_BCCut_File}.
%
The first line is a required header intended to assist the user,
and is otherwise ignored.
%
Subsequent lines define each named cut and the block information on
either side of the interface.
%
This information consists of the block number, face direction and plane,
and the subregion of that face.
%
To assist the user, each line is parsed based on whitespace rather than
assuming a fixed format.  
%
Comment lines beginning with ``\,!'' can be used throughout the file
and blank lines are ignored.
%
In order to perform the matching, each cut name must be unique.
%
Face and subregion index directions are specified as I, J, or K.
%
Face and subregion locations are specified as MIN, MAX, or as an
integer value.  
%
For two-dimensional flows, all K index values should be set to MIN.
%
Checks are performed to ensure that the specified interfaces are
complete and consistent.

%%%%%%%%%%%%%%%%%%%%%%%%%%%%%%%%%%%%%%%%%%%%%%%%%%%%%%%%%%%%%%%%%%%%%%%%%%%%%%%%

\vspace{1.0em}%
\IDEA{
\begin{itemize}
\setlength\topsep{0.0em}
\setlength\itemsep{0.0em}
\setlength\parsep{0.0em}
\item \raggedright
\Lref{LST:Plot3d_BCCut_File} might benefit from a more detailed example
with multiple cuts,
left\slash right cuts with differing number of subfaces, or
subregion locations that are not simply min\slash max.
One example is the 3-block 2d nozzle example from the Wind-US tutorial.
\item \raggedright
Need to emphasize each cut is a pair and the name matches between the two lines
that specify each half of the pair.
\end{itemize}
%
} %end FIXME

%%%%%%%%%%%%%%%%%%%%%%%%%%%%%%%%%%%%%%%%%%%%%%%%%%%%%%%%%%%%%%%%%%%%%%%%%%%%%%%%

\begin{center}
% Use minipage to force the lstlisting width. Else it will be \textwidth.
% Use this rule to visually determine needed width value.
%\rule{0.78\textwidth}{2.0pt}
\begin{minipage}{0.78\textwidth}
\begin{lstlisting}[language={},basicstyle=\small\ttfamily
,label={LST:Plot3d_BCCut_File}
,caption={Plot3D interface cut file example.},captionpos=b
]
CUT_NAME  BLK  FACE  PLACE  IND1  BEG  END  IND2  BEG  END
!=======  ===  ====  =====  ====  ===  ===  ====  ===  ===
ONE         1    I    MAX     J   MIN  MAX    K   MIN  MAX
ONE         2    I    MIN     J   MIN  MAX    K   MIN  MAX
\end{lstlisting}
%% Label must go in lstlisting optional args!
%% Anything inside the listing gets typeset!
%\label{LST:Plot3d_BCCut_File} % Keep this for grep to find.
\end{minipage}
\end{center}

%%%%%%%%%%%%%%%%%%%%%%%%%%%%%%%%%%%%%%%%%%%%%%%%%%%%%%%%%%%%%%%%%%%%%%%%%%%%%%%%

\SeeAlso{
\inp[nmlGrid::Plot3dCutfile]{Plot3d\_Cutfile}
}

%%%%%%%%%%%%%%%%%%%%%%%%%%%%%%%%%%%%%%%%%%%%%%%%%%%%%%%%%%%%%%%%%%%%%%%%%%%%%%%%


\paragraph{Plot3D Boundary Condition File}

%%%%%%%%%%%%%%%%%%%%%%%%%%%%%%%%%%%%%%%%%%%%%%%%%%%%%%%%%%%%%%%%%%%%%%%%%%%%%%%%

For Plot3D grids, boundary conditions are specified via a separate
text-based boundary condition file.
%
An example of this file is shown in \cref{LST:Plot3d_BCSfile}.
%
Once again, the first line is a header intended to assist the user,
and is otherwise ignored.
%
%The format of the header is very similar to that in the
%Plot3d interface cut file.
%
Subsequent lines define each named boundary condition and the
block information for where that condition applies.
%
This information consists of the block number, face direction and plane,
and the subregion of that face.
%
This information is parsed based on whitespace rather than assuming a
fixed format.
%
Comment lines beginning with ``\,!'' can be used throughout the
file and blank lines are ignored.
%
Face and subregion index directions are specified as I, J, or K.
%
Face and subregion locations are specified as MIN, MAX, or as an
integer value.
%
For two-dimensional flows, all K index values should be set to MIN.
%
The supported boundary condition names are listed in
\tref{TAB:Plot3dBCSFileBCNames}
and are used to specify the type of boundary condition to apply
to a particular face region.
%
\begin{table}
\begin{center}
\caption{Supported boundary condition names in Plot3D boundary condition file.}
\label{TAB:Plot3dBCSFileBCNames}
\begin{tabular}{ll}
  BC\_NAME              &  Description \\
  \hline
  INFLOW                &  Generic Inflow \\
  SUBINFLOW             &  Subsonic Inflow \\
  SUPINFLOW             &  Supersonic Inflow \\
  OUTFLOW               &  Generic Outflow \\
  SUBOUTFLOW            &  Subsonic Outflow \\
  SUPOUTFLOW \verb$|$
  EXTRAP                &  Supersonic Outflow (Extrapolation) \\
  FARFIELD              &  Generic farfield \\
  CHARACTERISTIC        &  Characteristic Boundary \\
  FREESTREAM            &  Freestream \\
  FIXED                 &  Fixed \\
  SLIP                  &  Inviscid Wall \\
  AWALL \verb$|$        
  WALL                  &  Adiabatic Viscous Wall \\
  IWALL                 &  Isothermal Viscous Wall \\
  SYMMETRY              &  Symmetry \\
  PERIODIC              &  Periodic
\end{tabular}
\end{center}
\end{table}
%
The group name is a user-specified string that is used to specify
a named set of flow values to apply in conjunction with the
boundary condition type.
%
The group name should correspond to a
\BCinpBC[nmlBC::BCInput/Name]{Name}
entry in the namelist input file.
%
Numerical values for the flow conditions will then be taken from
the other \BCinpBC[nmlBC::BCInput]{} properties.

%%%%%%%%%%%%%%%%%%%%%%%%%%%%%%%%%%%%%%%%%%%%%%%%%%%%%%%%%%%%%%%%%%%%%%%%%%%%%%%%

\begin{center}
% Use minipage to force the lstlisting width. Else it will be \textwidth.
% Use this rule to visually determine needed width value.
%\rule{0.95\textwidth}{2.0pt}
\begin{minipage}{0.95\textwidth}
\begin{lstlisting}[language={},basicstyle=\small\ttfamily
,label={LST:Plot3d_BCSfile}
,caption={Plot3D boundary conditions file example.},captionpos=b
]
BC_NAME   BLK  FACE  PLACE  IND1  BEG  END  IND2  BEG  END   GROUPNAME
!======   ===  ====  =====  ====  ===  ===  ====  ===  ===   =========
SUBINFLOW   1    I    MIN     J   MIN  MAX    K   MIN  MAX   INFLOW
AWALL       1    J    MIN     I   MIN  MAX    K   MIN  MAX   BOT_WALL
AWALL       1    J    MAX     I   MIN  MAX    K   MIN  MAX   TOP_WALL
!
EXTRAP      2    I    MAX     J   MIN  MAX    K   MIN  MAX   OUTFLOW
AWALL       2    J    MIN     I   MIN  MAX    K   MIN  MAX   BOT_WALL
AWALL       2    J    MAX     I   MIN  MAX    K   MIN  MAX   TOP_WALL
\end{lstlisting}
%% Label must go in lstlisting optional args!
%% Anything inside the listing gets typeset!
%\label={LST:Plot3d_BCSfile} % keep for grep to find
\end{minipage}
\end{center}

%%%%%%%%%%%%%%%%%%%%%%%%%%%%%%%%%%%%%%%%%%%%%%%%%%%%%%%%%%%%%%%%%%%%%%%%%%%%%%%%

\vspace{1.0em}%
\IDEA{
\begin{itemize}
\setlength\topsep{0.0em}
\setlength\itemsep{0.0em}
\setlength\parsep{0.0em}
\item \raggedright
Need better distinction between FARFIELD, CHARACTERISTIC, FREESTREAM.
\item \raggedright
\Lref{LST:Plot3d_BCSfile} might also benefit from using the same
3-block 2d nozzle example from the Wind-US tutorial.
\end{itemize}
%
} %end IDEA

%%%%%%%%%%%%%%%%%%%%%%%%%%%%%%%%%%%%%%%%%%%%%%%%%%%%%%%%%%%%%%%%%%%%%%%%%%%%%%%%

\SeeAlso{
\inp[nmlGrid::Plot3dBCSfile]{Plot3d\_BCSfile}
}

%%%%%%%%%%%%%%%%%%%%%%%%%%%%%%%%%%%%%%%%%%%%%%%%%%%%%%%%%%%%%%%%%%%%%%%%%%%%%%%%


\paragraph{Limitations}

%%%%%%%%%%%%%%%%%%%%%%%%%%%%%%%%%%%%%%%%%%%%%%%%%%%%%%%%%%%%%%%%%%%%%%%%%%%%%%%%

GFR does not support Plot3D iblanking.
%
This feature is often used to cut holes in structured grids and omit those
points from the flow domain, or to define donor interpolation information
for chimera overset meshes.
%
Depending on the nature of the iblanking, it may be possible to decompose the
mesh into multiple zones that no longer require iblanking.

%%%%%%%%%%%%%%%%%%%%%%%%%%%%%%%%%%%%%%%%%%%%%%%%%%%%%%%%%%%%%%%%%%%%%%%%%%%%%%%%

